\subsection{Truth Tables}
\subsubsection{Exercises:}
\begin{enumerate}
\item Make truth tables for the following formulas:
  \begin{enumerate} 
  \item $\neg P \lor Q$.
  \item $(S \lor G) \land (\neg S \lor \neg G)$.
  \end{enumerate}
  
\item Make truth tables for the following formulas:
  \begin{enumerate}
  \item $\neg [P \land (Q \lor \neg P)]$.
  \item $(P \lor Q) \land (\neg P \lor R)$.
  \end{enumerate}

\item In this exercise we will use the symbol $+$ to mean \emph{exclusive or}. In other words, $P + Q$ means ``$P$ or $Q$, but not both.''
  \begin{enumerate}
  \item Make a truth table for $P + Q$.
  \item Find a formula using only the connectives $\land$, $\lor$, and $\neg$ that is equivalent to $P + Q$. Justify your answer with a truth table.
  \end{enumerate}

\item Find a formula using only the connectives $\land$ and $\neg$ that is equivalent to $P \lor Q$. Justify your answer with a truth table.

\item Some mathematicians use the symbol $\downarrow$ to mean \emph{nor}. In other words $P \downarrow Q$ means ``neither P nor Q.''
  \begin{enumerate}
  \item Make a truth table for $P \downarrow Q$.
  \item Find a formula using only the connectives $\land$, $\lor$, and $\neg$ that is equivalent to $P \downarrow Q$.
  \item Find formulas using only the connective $\downarrow$ that are equivalent to $\neg P$, $P \lor Q$, and $P \land Q$.
  \end{enumerate}

\item Some mathematicians write $P \mid Q$ to mean ``$P$ and $Q$ are not both true.'' (This connective is called \emph{nand}, and is used in the study of circuits in computer science.)
  \begin{enumerate}
  \item Make a truth table for $P \mid Q$.
  \item Find a formula using only the connective $\land$, $\lor$, and $\neg$ that is equivalent to $P \mid Q$.
  \item Find formulas using only the connective $\mid$ that are equivalent to $\neg P$, $P \lor Q$, and $P \land Q$.
  \end{enumerate}

\item Use truth tables to determine whether or not the arguments in exercise $7$ of Section $1.1$ are valid.
  
\item Use truth tables to determine which of the following formulas are equivalent to each other:
  \begin{enumerate}
  \item $(P \land Q) \lor (\neg P \land \neg Q)$.
  \item $\neg P \land Q$.
  \item $(P \lor \neg Q) \land (Q \lor \neg P)$.
  \item $\neg (P \lor Q)$.
  \item $(Q \land P) \lor \neg P$.
  \end{enumerate}

\item Use truth tables to determine which of these statements are tautologies, which are contradictions, and which are neither:
  \begin{enumerate}
  \item $(P \lor Q) \land (\neg P \lor \neg Q)$.
  \item $(P \lor Q) \land (\neg P \land \neg Q)$.
  \item $(P \lor Q) \lor (\neg P \lor \neg Q)$.
  \item $[P \land (Q \lor \neg R)] \lor (\neg P \lor R)$.
  \end{enumerate}

\item Use truth tables to check these laws:
  \begin{enumerate}
  \item The second DeMorgan's law. (The first was checked in the text.)
  \item The distributive laws.
  \end{enumerate}

\item Use the laws stated in the text to find simpler formulas equivalent to these formulas. (See Examples $1.2.5$ and $1.2.7$.)
  \begin{enumerate}
  \item $\neg (\neg P \land \neg Q)$.
  \item $(P \land Q) \lor (P \land \neg Q)$.
  \item $\neg (P \land \neg Q) \lor (\neg P \land Q)$.
  \end{enumerate}

\item Use the laws stated in the text to find simpler formulas equivalent to these formulas. (See examples $1.2.5$ and $1.2.7$.)
  \begin{enumerate}
  \item $\neg (\neg P \lor Q) \lor (P \land \neg R)$.
  \item $\neg (\neg P \land Q) \lor (P \lor \neg R)$.
  \item $(P \land R) \lor [\neg R \land (P \lor Q)]$.
  \end{enumerate}

\item Use the first DeMorgan's law and the double negation law to derive the second DeMorgan's law.

\item Note that the associative laws say only that parentheses are unecessary when combining \emph{three} statements with $\land$ or $\lor$. In fact, these laws can be used to justify leaving parentheses out when more than three statements are combined. Use associative laws to show that $[P \land (Q \land R)] \land S$ is equivalent to $(P \land Q) \land (R \land S)$.

\item How many lines will there be in the truth table for a statement containing $n$ letters?

\item Find a formula involving the connectives $\land$, $\lor$, and $\neg$ that has the following truth table:
  \begin{center}
    \begin{tabular}{ccc}
      $P$ & $Q$ & $???$ \\
      \midrule
      F & F & T \\
      F & T & F \\
      T & F & T \\
      T & T & T
    \end{tabular}
  \end{center}

\item Find a formula involving the connectives $\land$, $\lor$, and $\neg$ that has the following truth table:
  \begin{center}
    \begin{tabular}{ccc}
      $P$ & $Q$ & $???$ \\
      \midrule
      F & F & F \\
      F & T & T \\
      T & F & T \\
      T & T & F
    \end{tabular}
  \end{center}

\item Suppose the conclusion of an argument is a tautology. What can you conclude about the validity of the argument? What if the conclusion is a contradiction? What if one of the premises is either a tautology or a contradiction?
\end{enumerate}

\subsubsection{Solutions:}
\begin{enumerate}
\item Truth Tables:
  \begin{enumerate}
  \item \textbf{Formula:} $\neg P \lor Q$. \\
    \textbf{Truth Table:}
    \begin{center}
      \begin{tabular}{cccc}
        $P$ & $Q$ & $\neg P$ & $\neg P \lor Q$ \\
        \midrule
        T & T & F & T \\
        T & F & F & F \\
        F & T & T & T \\
        F & F & T & T
      \end{tabular}
    \end{center}
    
  \item \textbf{Formula:} $(S \lor G) \land (\neg S \lor \neg G)$. \\
    \textbf{Truth Table:}
    \begin{center}
      \begin{tabular}{ccccc}
        $S$ & $G$ & $S \lor G$ & $\neg S \lor \neg G$ & $(S \lor G) \land (\neg S \lor \neg G)$ \\
        \midrule
        T & T & T & F & F \\
        T & F & T & T & T \\
        F & T & T & T & T \\
        F & F & F & T & F 
      \end{tabular}
    \end{center}
  \end{enumerate}

\item Truth Tables:
  \begin{enumerate}
  \item \textbf{Formula:} $\neg [P \land (Q \lor \neg P)]$. \\
    \textbf{Truth Table:}
    \begin{center}
      \begin{tabular}{cccccc}
        $P$ & $Q$ & $\neg P$ & $Q \lor \neg P$ & $P \land (Q \lor \neg P)$ & $\neg[P \land (Q \lor \neg P)]$ \\
        \midrule
        T & T & F & T & T & F \\ 
        T & F & F & F & F & T \\
        F & T & T & T & F & T \\ 
        F & F & T & T & F & T 
      \end{tabular}
    \end{center}

  \item \textbf{Formula:} $(P \lor Q) \land (\neg P \lor R)$. \\
    \textbf{Truth Table:}
    \begin{center}
      \begin{tabular}{cccccc}
        $P$ & $Q$ & $R$ & $P \lor Q$ & $\neg P \lor R$ & $(P \lor Q) \land (\neg P \lor R)$ \\
        \midrule
        T & T & T & T & T & T \\
        T & T & F & T & F & F \\
        T & F & T & T & T & T \\
        T & F & F & T & F & F \\
        F & T & T & T & T & T \\ 
        F & T & F & T & T & T \\
        F & F & T & F & T & F \\
        F & F & F & F & T & F
      \end{tabular}
    \end{center}
  \end{enumerate}

\item Exclusive Or ($+$):
  \begin{enumerate}
  \item \textbf{Truth Table for } $P + Q$:
    \begin{center}
      \begin{tabular}{ccc}
        $P$ & $Q$ & $P + Q$ \\
        \midrule
        T & T & F \\
        T & F & T \\
        F & T & T \\
        F & F & F
      \end{tabular}
    \end{center}

  \item \textbf{Equivalent Formula:} $(P \lor Q) \land \neg (P \land Q)$. \\
    \textbf{Reasoning:} ``$P$ or $Q$'' translates to $(P \lor Q)$, and ``but not both'' translates to $\land \neg (P \land Q)$.
  \end{enumerate}

\item Equivalent to $P \lor Q$ Using Only $\land$ and $\neg$: \\
  \textbf{Formula:} $\neg (\neg P \land \neg Q)$. \\
  \textbf{Truth Table}:
  \begin{center}
    \begin{tabular}{cccc}
      $P$ & $Q$ & $P \lor Q$ & $\neg (\neg P \land \neg Q)$ \\
      \midrule
      T & T & T & T \\
      T & F & T & T \\
      F & T & T & T \\
      F & F & F & F 
    \end{tabular}
  \end{center}
  \textbf{Reasoning:} By DeMorgan's laws, negating the conjunction of two negated variables yields the disjunction of the original variables.

\item Nor $(P \downarrow Q)$:
  \begin{enumerate}
  \item \textbf{Truth Table for} $P \downarrow Q$:
    \begin{center}
      \begin{tabular}{ccc}
        $P$ & $Q$ & $P \downarrow Q$ \\
        \midrule
        T & T & F \\
        T & F & F \\
        F & T & F \\
        F & F & T 
      \end{tabular}
    \end{center}
  \item \textbf{Equivalent Formula:} $\neg (P \lor Q)$.
  \item \textbf{Equivalent Formulas Using Only $\downarrow$}:
    \begin{itemize}
    \item $\neg P \equiv P \downarrow P$.
    \item $P \lor Q \equiv (P \downarrow Q) \downarrow (P \downarrow Q)$.
    \item $P \land Q \equiv (P \downarrow P) \downarrow (Q \downarrow Q)$.
    \end{itemize}
    \textbf{Reasoning:} ``Neither $P$ nor $Q$'' is the exact negation of ``$P$ or $Q$''.
  \end{enumerate}

\item Nand $(P \mid Q)$:
  \begin{enumerate}
  \item \textbf{Truth Table for} $P \mid Q$:
    \begin{center}
      \begin{tabular}{ccc}
        $P$ & $Q$ & $P \mid Q$ \\
        \midrule
        T & T & F \\
        T & F & T \\
        F & T & T \\
        F & F & T 
      \end{tabular}
    \end{center}

  \item \textbf{Equivalent Formula:} $\neg (P \land Q)$.
    
  \item \textbf{Equivalent Formulas Using Only $\mid$:}
    \begin{itemize}
    \item $\neg P \equiv P \mid P$.
    \item $P \lor Q \equiv (P \mid P) \mid (Q \mid Q)$.
    \item $P \land Q \equiv (P \mid Q) \mid (P \mid Q)$.
    \end{itemize}
    \textbf{Reasoning:} ``Not both true'' is the direct negation of the conjunction.
  \end{enumerate}

\item Validity Analysis via Truth Tables:
  \begin{enumerate}
  \item Argument (a):
    \textbf{Premises:} $\neg (J_M \land P_M), P_M \lor P_C, J_M$. \\
    \textbf{Conclusion:} $P_C$. \\
    \textbf{Truth Table:}
    \begin{center}
      \begin{tabular}{cccccccc}
        $J_M$ & $P_M$ & $P_C$ & $J_M \land P_M$ & $\neg (J_M \land P_M)$ & $P_M \lor P_C$ & $J_M$ & $P_C$ \\
        \midrule
        T & T & T & T & F & T & T & T \\
        T & T & F & T & F & T & T & F \\
        T & F & T & F & T & T & T & T \\
        T & F & F & F & T & F & T & F \\ 
        F & T & T & F & T & T & F & T \\
        F & T & F & F & T & T & F & F \\
        F & F & T & F & T & T & F & T \\
        F & F & F & F & T & F & F & F
      \end{tabular}
    \end{center}
    \textbf{Result:} \emph{Valid}. In the only row where all three premises are True (Row 3). the conclusion $(P_C)$ is also True.

  \item Argument (b): \\
    \textbf{Premises:} $B \lor F, P \lor C, \neg(F \land C)$. \\
    \textbf{Conclusion:} $\neg(B \land P)$. \\
    \textbf{Truth Table:}
    \begin{center}
      \begin{tabular}{cccccccc}
        $B$ & $F$ & $P$ & $C$ & $B \lor F$ & $P \lor C$ & $\neg(F \land C)$ & $\neg(B \land P)$ \\
        \midrule
        T & T & T & T & T & T & F & F \\
        T & T & T & F & T & T & T & F \\
        T & T & F & T & T & T & F & T \\
        T & T & F & F & T & F & T & T \\
        T & F & T & T & T & T & T & F \\
        T & F & T & F & T & T & T & F \\
        T & F & F & T & T & T & T & T \\
        T & F & F & F & T & F & T & T \\
        F & T & T & T & T & T & F & T \\
        F & T & T & F & T & T & T & T \\
        F & T & F & T & T & T & F & T \\
        F & T & F & F & T & F & T & T \\
        F & F & T & T & F & T & T & T \\
        F & F & T & F & F & T & T & T \\
        F & F & F & T & F & T & T & T \\
        F & F & F & F & F & F & T & T 
      \end{tabular}
    \end{center}
    \textbf{Result:} \emph{Invalid}. In the 2nd, 5th, and 6th rows, all three premises are True, but the conclusion is False.

  \item Argument (c): \\
    \textbf{Premises:} $J \lor B, \neg S \lor \neg B$. \\
    \textbf{Conclusion:} $J \lor \neg S$. \\
    \textbf{Truth Table:}
    \begin{center}
      \begin{tabular}{cccccccc}
        $J$ & $B$ & $S$ & $\neg B$ & $\neg S$ & $J \lor B$ & $\neg S \lor \neg B$ & $J \lor \neg S$ \\
        \midrule
        T & T & T & F & F & T & F & T \\
        T & T & F & F & T & T & T & T \\
        T & F & T & T & F & T & T & T \\
        T & F & F & T & T & T & T & T \\
        F & T & T & F & F & T & F & F \\
        F & T & F & F & T & T & T & T \\
        F & F & T & T & F & F & T & F \\
        F & F & F & T & T & F & T & T
      \end{tabular}
    \end{center}
    \textbf{Result:} \emph{Valid}. In every row where both premises are True (Rows 2, 3, 4, and 6), the conclusion is also True.

  \item Argument (d): \\
    \textbf{Premises:} $(S \land H) \lor (E \land \neg H)$. \\
    \textbf{Conclusion:} $\neg(S \land E)$. \\
    \textbf{Truth Table:}
    \begin{center}
      \begin{tabular}{cccccccc}
        $S$ & $E$ & $H$ & $S \land H$ & $E \land \neg H$ & $(S \land H) \lor (E \land \neg H)$ & $S \land E$ & $\neg(S \land E)$ \\
        \midrule
        T & T & T & T & F & T & T & F \\
        T & T & F & F & T & T & T & F \\
        T & F & T & T & F & T & F & T \\
        T & F & F & F & F & F & F & T \\
        F & T & T & F & F & F & F & T \\
        F & T & F & F & T & T & F & T \\
        F & F & T & F & F & F & F & T \\
        F & F & F & F & F & F & F & T
      \end{tabular}
    \end{center}
    \textbf{Result:} \emph{Invalid}. In the 1st and 2nd rows, the premise is True, but the conclusion is False.
  \end{enumerate}

\item Equivalence Analysis via Truth Tables: \\
  \textbf{Truth Table:}
  \begin{center}
    \begin{tabular}{ccccccc}
      $P$ & $Q$ & $(P \land Q) \lor (\neg P \land \neg Q)$ & $\neg P \land Q$ & $(P \lor \neg Q) \land (Q \lor \neg P)$ & $\neg (P \lor Q)$ & $(Q \land P) \lor \neg P$ \\
      \midrule
      T & T & T & F & T & F & T \\
      T & F & F & F & F & F & F \\
      F & T & F & T & F & F & T \\
      F & F & T & F & T & T & T 
    \end{tabular}
  \end{center}
  \textbf{Result:} The colums $(P \land Q) \lor (\neg P \land \neg Q)$ and $(P \lor \neg Q) \land (Q \lor \neg P)$ have the exact same truth values in every row, Therefore $(P \land Q) \lor (\neg P \land \neg Q) \equiv (P \lor \neg Q) \land (Q \lor \neg P)$.

\item Tautology, Contradiction, or Neither:
  \begin{enumerate}
  \item[(a)] $(P \lor Q) \land (\neg P \lor \neg Q)$
  \item[(b)] $(P \lor Q) \land (\neg P \land \neg Q)$
  \item[(c)] $(P \lor Q) \lor (\neg P \lor \neg Q)$
  \item[(d)] $[P \land (Q \lor \neg R)] \lor (\neg P \lor R)$
  \end{enumerate}
  \textbf{Truth Table:}
  \begin{center}
    \begin{tabular}{ccccccc}
      $P$ & $Q$ & $R$ & (a) & (b) & (c) & (d) \\
      \midrule
      T & T & T & F & F & T & T \\
      T & T & F & F & F & T & T \\
      T & F & T & T & F & T & T \\
      T & F & F & T & F & T & T \\
      F & T & T & T & F & T & T \\
      F & T & F & T & F & T & T \\
      F & F & T & F & F & T & T \\
      F & F & F & F & F & T & T
    \end{tabular}
  \end{center}
  \textbf{Result:}
  \begin{itemize}
  \item (a) is \emph{Neither} (True in some rows, False in others).
  \item (b) is a \emph{Contradiction} (False in all rows).
  \item (c) is a \emph{Tautology} (True in all rows).
  \item (d) is a \emph{Tautology} (True in all rows).
  \end{itemize}

\item Checking Laws:
  \begin{enumerate}
  \item \textbf{The second DeMorgan's law:} ($\neg(P \land Q) \equiv \neg P \lor \neg Q$)
    \begin{center}
      \begin{tabular}{ccccccc}
        $P$ & $Q$ & $P \land Q$ & $\neg(P \land Q)$ & $\neg P$ & $\neg Q$ & $\neg P \lor \neg Q$ \\
        \midrule
        T & T & T & \textbf{F} & F & F & \textbf{F} \\
        T & F & F & \textbf{T} & F & T & \textbf{T} \\
        F & T & F & \textbf{T} & T & F & \textbf{T} \\
        F & F & F & \textbf{T} & T & T & \textbf{T}
      \end{tabular}
    \end{center}
    \textbf{Result:} The column values for $\neg(P \land Q)$ and $\neg P \lor \neg Q$ match perfectly, proving the law.

    \item \textbf{The Distributive Laws} \\
    \textit{Law 1:} $P \land (Q \lor R) \equiv (P \land Q) \lor (P \land R)$
    \begin{center}
      \begin{tabular}{cccccccc}
        $P$ & $Q$ & $R$ & $Q \lor R$ & $P \land (Q \lor R)$ & $P \land Q$ & $P \land R$ & $(P \land Q) \lor (P \land R)$ \\
        \midrule
        T & T & T & T & \textbf{T} & T & T & \textbf{T} \\
        T & T & F & T & \textbf{T} & T & F & \textbf{T} \\
        T & F & T & T & \textbf{T} & F & T & \textbf{T} \\
        T & F & F & F & \textbf{F} & F & F & \textbf{F} \\
        F & T & T & T & \textbf{F} & F & F & \textbf{F} \\
        F & T & F & T & \textbf{F} & F & F & \textbf{F} \\
        F & F & T & T & \textbf{F} & F & F & \textbf{F} \\
        F & F & F & F & \textbf{F} & F & F & \textbf{F}
      \end{tabular}
    \end{center}
    \textbf{Result:} The values match perfectly, proving the first distributive law.

    \textit{Law 2:} $P \lor (Q \land R) \equiv (P \lor Q) \land (P \lor R)$
    \begin{center}
      \begin{tabular}{cccccccc}
        $P$ & $Q$ & $R$ & $Q \land R$ & $P \lor (Q \land R)$ & $P \lor Q$ & $P \lor R$ & $(P \lor Q) \land (P \lor R)$ \\
        \midrule
        T & T & T & T & \textbf{T} & T & T & \textbf{T} \\
        T & T & F & F & \textbf{T} & T & T & \textbf{T} \\
        T & F & T & F & \textbf{T} & T & T & \textbf{T} \\
        T & F & F & F & \textbf{T} & T & T & \textbf{T} \\
        F & T & T & T & \textbf{T} & T & T & \textbf{T} \\
        F & T & F & F & \textbf{F} & T & F & \textbf{F} \\
        F & F & T & F & \textbf{F} & F & T & \textbf{F} \\
        F & F & F & F & \textbf{F} & F & F & \textbf{F}
      \end{tabular}
    \end{center}
    \textbf{Result:} The values match perfectly, proving the second distributive law.
  \end{enumerate}

\item Simpler Formulas:
  \begin{enumerate}
  \item $\neg (\neg P \land \neg Q)$ 
    \begin{align*}
      \neg (\neg P \land \neg Q) &\equiv \neg \neg P \lor \neg \neg Q \quad \text{DeMorgan's law} \\ 
      &\equiv P \lor Q \quad \text{double negation law}
    \end{align*}

  \item $(P \land Q) \lor (P \land \neg Q)$ \\
    \begin{align*}
      (P \land Q) \lor (P \land \neg Q) &\equiv P \land (Q \lor \neg Q) \quad \text{Distributive law} \\
      &\equiv P \land True \quad \text{tautology law} \\
      &\equiv P \quad \text{identity law}
    \end{align*}

  \item $\neg (P \land \neg Q) \lor (\neg P \land Q)$ \\
    \begin{align*}
      \neg (P \land \neg Q) \lor (\neg P \land Q) &\equiv (\neg P \lor \neg \neg Q) \lor (\neg P \land Q) \quad \text{DeMorgan's law} \\
      &\equiv (\neg P \lor Q) \lor (\neg P \land Q) \quad \text{double negation law} \\
      &\equiv \neg P \lor (Q \lor (\neg P \land Q)) \quad \text{associative law} \\
      &\equiv \neg P \lor ((Q \lor \neg P) \land (Q \lor Q)) \quad \text{distributive law} \\
      &\equiv \neg P \lor ((Q \lor \neg P) \land Q) \quad \text{idempotent law} \\
      &\equiv \neg P \lor Q \quad \text{absorption law}
    \end{align*}
  \end{enumerate}

\item Simpler Formulas:
  \begin{enumerate}
  \item $\neg (\neg P \lor Q) \lor (P \land \neg R)$ \\
    \begin{align*}
      \neg (\neg P \lor Q) \lor (P \land \neg R) &\equiv (\neg \neg P \land \neg Q) \lor (P \lor \neg R) \quad \text{DeMorgan's law} \\
      &\equiv (P \land \neg Q) \lor (P \land \neg R) \quad \text{double negation law} \\
      &\equiv P \land (\neg Q \lor \neg R) \quad \text{distributive law} \\
      &\equiv P \land \neg(Q \land R) \quad \text{DeMorgan's law}
    \end{align*}

  \item $\neg (\neg P \land Q) \lor (P \lor \neg R)$ \\
    \begin{align*}
      \neg (\neg P \land Q) \lor (P \lor \neg R) &\equiv (\neg \neg P \lor \neg Q) \lor (P \lor \neg R) \quad \text{DeMorgan's law} \\
      &\equiv (P \lor \neg Q) \lor (P \lor \neg R) \quad \text{double negation law} \\
      &\equiv (P \lor P) \lor (\neg Q \lor \neg R) \quad \text{associative and commutative laws} \\
      &\equiv P \lor (\neg Q \lor \neg R) \quad \text{idempotent law} \\ 
      &\equiv P \lor \neg (Q \land R) \quad \text{DeMorgan's law}
    \end{align*}

  \item $(P \land R) \lor [\neg R \land (P \lor Q)]$ \\
    \begin{align*}
      (P \land R) \lor [\neg R \land (P \lor Q) &\equiv (P \land R) \lor [(\neg R \land P) \lor (\neg R \land Q)] \quad \text{distributive law} \\
        &\equiv [(P \land R) \lor (\neg R \land P)] \lor (\neg R \land Q) \quad \text{associative law} \\
        &\equiv [(R \land P) \lor (\neg R \land P)] \lor (\neg R \land Q) \quad \text{commutative law} \\
        &\equiv [(R \lor \neg R) \land P] \lor (\neg R \land Q) \quad \text{distributive law} \\
        &\equiv [True \land P] \lor (\neg R \land Q) \quad \text{tautology law} \\
        &\equiv P \lor (\neg R \land Q) \quad \text{identity law}
    \end{align*}
  \end{enumerate}

  \item \textbf{Deriving the Second DeMorgan's Law:} \\
  Start with the negation of a conjunction, $\neg(P \land Q)$. \\
  By the double negation law ($P \equiv \neg\neg P$), we can rewrite the inner variables $P$ and $Q$:
  \[ \neg(P \land Q) \equiv \neg(\neg\neg P \land \neg\neg Q) \]
  Applying the first DeMorgan's law, $\neg(A \land B) \equiv \neg A \lor \neg B$, to the terms inside the outer negation:
  \[ \neg(\neg\neg P \land \neg\neg Q) \equiv \neg\neg(\neg P \lor \neg Q) \]
  Finally, applying the double negation law again removes the two outermost negations:
  \[ \neg\neg(\neg P \lor \neg Q) \equiv \neg P \lor \neg Q \]
  Therefore, $\neg(P \land Q) \equiv \neg P \lor \neg Q$.

  \item \textbf{Associative Laws for Multiple Statements:} \\
  By the associative law for conjunction, $A \land (B \land C) \equiv (A \land B) \land C$. \\
  Applying this to the inner expression with $A=P$, $B=Q$, and $C=R$:
  \[ [P \land (Q \land R)] \land S \equiv [(P \land Q) \land R] \land S \]
  Now, applying the associative law a second time treating $A=(P \land Q)$, $B=R$, and $C=S$:
  \[ [(P \land Q) \land R] \land S \equiv (P \land Q) \land (R \land S) \]

  \item \textbf{Truth Table Lines:} \\
  There will be $2^n$ lines. Since each distinct letter has exactly 2 possible truth values (True or False), the total number of combinations for $n$ variables is $2 \times 2 \times \dots \times 2$ ($n$ times).

  \item \textbf{Formula from Truth Table:}
  The formula evaluates to False in exactly one case: when $P$ is False and $Q$ is True. In all other cases, it is True. \\
  A logical statement that is true whenever $P$ is true OR whenever $Q$ is false matches this perfectly. \\
  Therefore, the equivalent formula is: $P \lor \neg Q$.

  \item \textbf{Formula from Truth Table:}
  This truth table represents the exclusive OR (XOR) operation, which evaluates to True if and only if exactly one of $P$ or $Q$ is True. \\
  Using standard connectives, this means "$P$ or $Q$, but not both $P$ and $Q$". \\
  The equivalent formula is: $(P \lor Q) \land \neg(P \land Q)$.

  \item \textbf{Analyzing Argument Conclusions and Premises:} \\
  \begin{itemize}
    \item \textbf{If the conclusion is a tautology:} The argument is unconditionally valid. An argument is only invalid if it's possible for the premises to be true while the conclusion is false. Since a tautology is always true, this scenario is impossible.
    \item \textbf{If the conclusion is a contradiction:} The argument is valid if and only if the conjunction of all premises is also a contradiction. Otherwise, there would be a case where all premises are true but the conclusion is false (since a contradiction is always false).
    \item \textbf{If one of the premises is a tautology:} The tautology is always true and thus imposes no restrictions on the truth values. The validity of the argument depends entirely on the remaining premises.
    \item \textbf{If one of the premises is a contradiction:} The argument is vacuously valid. Because a contradiction is always false, it is impossible for all premises to be true simultaneously. Hence, you can never have the invalidating condition of "true premises and a false conclusion."
  \end{itemize}
\end{enumerate}
